% 核心结论：作业六过程：异频叠加后的波形
% 图原型：纯几何／对比；按本题相位与电路拓扑绘制。
\documentclass[tikz,border=7pt]{standalone}
\usepackage{fontspec,xeCJK,amsmath,amssymb}
\IfFontExistsTF{Times New Roman}{\setmainfont{Times New Roman}}{\setmainfont{TeX Gyre Termes}}
\IfFontExistsTF{Heiti SC}{\setCJKmainfont[AutoFakeBold=1.5]{Heiti SC}}{\setCJKmainfont[AutoFakeBold=1.5]{FandolHei-Regular}}
\usetikzlibrary{arrows.meta,calc,angles,quotes}
\definecolor{ink}{HTML}{243745}
\definecolor{blue}{HTML}{275DAB}
\definecolor{green}{HTML}{177C63}
\definecolor{red}{HTML}{BD4935}
\definecolor{purple}{HTML}{715393}
\begin{document}
\begin{tikzpicture}[font=\fontsize{9.5}{12}\selectfont,text=ink,
 v/.style={line width=1.1pt,-{Stealth[length=2.4mm,width=1.5mm]}},
 axis/.style={ink!65,line width=.45pt,-{Stealth[length=1.7mm,width=1.1mm]}},
 aux/.style={ink!40,densely dashed,line width=.5pt},
 lab/.style={fill=white,inner sep=2pt},
 card/.style={draw=blue!25,fill=blue!3,rounded corners=3pt,text width=6cm,align=left,inner sep=10pt}]
\path[use as bounding box] (-8,-3.35) rectangle (8,3.35);
\node[font=\fontsize{11}{14}\selectfont\bfseries] at (0.00000,3.04000) {作业六过程：异频叠加后的波形};
\draw[axis] (-6.70000,0.00000) -- (6.80000,0.00000);
\draw[axis] (-6.70000,-2.10000) -- (-6.70000,2.30000);
\node[] at (-6.40000,2.55000) {$\text{瞬时值}/\mathrm{V}$};
\node[] at (6.70000,-0.35000) {$x=\omega_1 t$};
\draw[ink] (-6.70000,-0.07000) -- (-6.70000,0.07000);
\node[] at (-6.70000,-2.35000) {$0$};
\draw[ink] (-3.45000,-0.07000) -- (-3.45000,0.07000);
\node[] at (-3.45000,-2.35000) {$\pi/2$};
\draw[ink] (-0.20000,-0.07000) -- (-0.20000,0.07000);
\node[] at (-0.20000,-2.35000) {$\pi$};
\draw[ink] (3.05000,-0.07000) -- (3.05000,0.07000);
\node[] at (3.05000,-2.35000) {$3\pi/2$};
\draw[ink] (6.30000,-0.07000) -- (6.30000,0.07000);
\node[] at (6.30000,-2.35000) {$2\pi$};
\draw[blue,thick] plot[domain=0:2*pi,samples=240] ({-6.7+13*\x/(2*pi)},{0.09848987309384055*(6*sqrt(2)*cos(deg(\x)))});
\draw[green,thick] plot[domain=0:2*pi,samples=240] ({-6.7+13*\x/(2*pi)},{0.09848987309384055*(8*sqrt(2)*cos(2*deg(\x)+30))});
\draw[red,thick] plot[domain=0:2*pi,samples=240] ({-6.7+13*\x/(2*pi)},{0.09848987309384055*((6*sqrt(2)*cos(deg(\x)))+(8*sqrt(2)*cos(2*deg(\x)+30)))});
\node[] at (-7.08000,-1.00000) {$-10.2$};
\node[] at (-7.08000,1.00000) {$10.2$};
\node[lab,text=blue] at (-3.60000,2.25000) {$u_1$};
\node[lab,text=green] at (0.00000,2.25000) {$u_2$};
\node[lab,text=red] at (3.60000,2.25000) {$u$};
\node[text width=15cm,align=center] at (0.00000,-3.00000) {各曲线使用相同纵轴尺度；红线在每个时刻等于蓝线与绿线的代数和。};
\end{tikzpicture}
\end{document}
